\section{Experiments and Results}
\label{ch:experiments}

This chapter presents the experimental setup, the environments used for evaluation, and the results obtained. We specifically focus on validating the hypothesis that a Meta-Controller LLM can effectively resolve exploration bottlenecks in a sparse Multi-Agent Reinforcement Learning (MARL) environment through dynamic reward shaping, and subsequently, how Hierarchical Reinforcement Learning (HRL) further improves these capabilities.

\subsection{Experimental Setup}
\label{sec:exp_setup}

To rigorously evaluate the proposed architectures, we compare distinct configurations within our multi-agent resource-gathering environment, structured into two parts. 

For the \textbf{Dynamic Reward Shaping} experiments, we evaluate:
\begin{enumerate}
    \item \textbf{Sparse Baseline:} A standard MAPPO implementation relying purely on sparse environmental rewards (+1 for completing milestones).
    \item \textbf{LLM Dynamic (Gemini 3.1 Flash Lite):} The proposed ``Mixing Board'' architecture, utilizing Google's Gemini 3.1 Flash Lite model via API as the Meta-Controller to dynamically shift potential-based reward weights. 
    \item \textbf{LLM Dynamic (Local Qwen 2.5 7B):} An equivalent dynamic architecture utilizing a localized, smaller parameter model (Qwen 2.5 7B) to observe the efficacy of smaller, self-hosted LLMs. 
\end{enumerate}

For the \textbf{Hierarchical Reinforcement Learning (HRL)} experiments, we evaluate three configurations:
\begin{enumerate}
    \item \textbf{No-LoRA Base (Qwen 2.5 7B):} The un-finetuned base model serving as the HRL Meta-Controller.
    \item \textbf{QLoRA with Reflection (Qwen 2.5 7B):} A QLoRA fine-tuned model with counterfactual reflection triggered on option staleness.
    \item \textbf{QLoRA without Reflection (Qwen 2.5 7B):} The same fine-tuned model, with the reflection mechanism disabled.
\end{enumerate}

All runs were executed with identical MAPPO hyperparameters (clipping $\epsilon=0.2$, GAE $\lambda=0.95$, $\gamma=0.99$, across 128 parallel environments).

% ==========================================
% PART 1: DYNAMIC REWARD SHAPING
% ==========================================
\subsection{Part 1: Dynamic Reward Shaping}
\label{sec:exp_dynamic}

The most crucial metric validating the LLM Dynamic architecture is the progression through the environment's strict dependency Directed Acyclic Graph (DAG). 

\begin{figure}[H]
    \centering
    \includegraphics[width=\textwidth]{../../plots/baseline_vs_llmdynamic_results/subtasks_comparison.png}
    \caption{Subtask Completion Rates: Comparison between the Sparse Baseline and the LLM Dynamic models over 2000 training updates.}
    \label{fig:subtasks_comparison_dynamic}
\end{figure}

As shown in \cref{fig:subtasks_comparison_dynamic}, all three configurations conquer the early-game tasks (Wood, Stone) almost immediately. The Pickaxe subtask follows shortly after, with the Baseline reaching it at step~1.1M and Gemini at step~720K.

However, the Baseline completely collapses at the \emph{Iron} bottleneck. Without any dynamic reward shaping to guide exploration, the Baseline agents never discover how to consistently mine Iron (0\% final rate), resulting in a total failure to learn all downstream tasks (Sword, Armor, Bridge). 

In contrast, both LLM-guided runs break through the Iron bottleneck, but at remarkably different speeds. \textbf{Qwen} unlocks Iron at step~10M, subsequently discovering Sword and Armor at essentially the same step (10.1M). This rapid cascade demonstrates the snowball effect: once a critical bottleneck is broken, downstream tasks are unlocked quickly. \textbf{Gemini} unlocks Iron much later, at step~40.7M (approximately 4$\times$ slower than Qwen). However, once unlocked, the same cascade occurs: Sword and Armor appear at step~41.3M.

Neither run achieved Enemy defeats (0\% for all three configurations), indicating that the combat subtask represents a qualitatively different challenge beyond reward shaping---likely requiring specialised action coordination that the current shaping signal alone cannot provide.

\subsubsection{LLM Interventions in Action}
\label{sec:exp_interventions}

To understand exactly \emph{how} the LLMs broke the bottleneck, we analyse the evolution of the reward weights. The weights shown in \cref{fig:llm_weights_evolution} are the \emph{post-softmax} weights actually applied to the potential-based shaping.

\begin{figure}[H]
    \centering
    \begin{subfigure}{0.8\textwidth}
        \includegraphics[width=\textwidth]{../../plots/baseline_vs_llmdynamic_results/llm_weights_evolution_gemini.png}
        \caption{Gemini 3.1 Flash Lite Meta-Controller (198 interventions)}
    \end{subfigure}
    \vskip\baselineskip
    \begin{subfigure}{0.8\textwidth}
        \includegraphics[width=\textwidth]{../../plots/baseline_vs_llmdynamic_results/llm_weights_evolution_qwen.png}
        \caption{Local Qwen 2.5 7B Meta-Controller (20 interventions)}
    \end{subfigure}
    \caption{Evolution of the Dynamic Reward Weights (post-softmax, as applied to the PBRS signal).}
    \label{fig:llm_weights_evolution}
\end{figure}

\Cref{fig:llm_weights_evolution} reveals strikingly different weight allocation strategies. \textbf{Gemini's strategy} is characterised by \emph{broad amplification} of the early-chain tasks. $w_{\text{wood}}$ and $w_{\text{stone}}$ rapidly climb and stay high, while $w_{\text{iron}}$ remains near zero after the softmax (mean~0.01). This pattern reveals a counter-intuitive insight: despite the raw LLM outputs assigning $w_{\text{iron}} = 1.0$, the \emph{softmax normalisation} compresses this into a negligible effective weight because the LLM simultaneously assigns high values to wood, stone, and bridge.

\textbf{Qwen's strategy} exhibits a \emph{focused suppression-and-amplification} profile. After an initial exploration phase, $w_{\text{wood}}$ and $w_{\text{stone}}$ are aggressively dropped, while $w_{\text{iron}}$ rises sharply to a peak of 3.87 during the mid-training bottleneck. By zeroing out already-mastered tasks and concentrating shaping energy on the bottleneck, Qwen breaks through Iron 30M steps earlier than Gemini.

\subsubsection{Reward Analysis}
\label{sec:exp_rewards}

\begin{figure}[H]
    \centering
    \includegraphics[width=0.7\textwidth]{../../plots/baseline_vs_llmdynamic_results/reward_comparison.png}
    \caption{Average Environment Reward (True Milestones Reached).}
    \label{fig:env_rewards_dynamic}
\end{figure}

\Cref{fig:env_rewards_dynamic} presents the true average environment reward. Both LLM Dynamic runs achieve significantly higher returns than the Baseline, reaching an identical peak reward ($\approx$0.076). This suggests that once the bottleneck is broken, the final quality of the learned policy is comparable; the difference lies purely in the \emph{sample efficiency} driven by Qwen's highly focused interventions.

% ==========================================
% PART 2: HIERARCHICAL REINFORCEMENT LEARNING
% ==========================================
\subsection{Part 2: Hierarchical Reinforcement Learning}
\label{sec:exp_hrl}

While the dynamic shaping architecture successfully guided agents past initial bottlenecks (like Iron), it ultimately stalled at the combat phase (Enemy, Gold). Modifying the continuous reward landscape was insufficient to teach the agents the discrete, long-horizon macro-actions required for coordinated combat.

This limitation motivated the shift to a full \textbf{Hierarchical Reinforcement Learning (HRL)} architecture, where the LLM does not merely shape the rewards, but explicitly dictates discrete \emph{Options} (e.g., ``Craft Sword'', ``Defeat Enemy'') to the lower-level MAPPO policy. Three configurations were evaluated over 2000 optimization updates ($\approx$65.5M environment steps each):

\begin{enumerate}
    \item \textbf{No-LoRA Base:} The un-finetuned Qwen 2.5 7B model serving as Meta-Controller.
    \item \textbf{QLoRA (With Reflection):} The QLoRA fine-tuned model with counterfactual reflection logic, triggered when an assigned Option becomes stale (agents fail to complete it within 150 steps).
    \item \textbf{QLoRA (No Reflection):} The identical fine-tuned model, but with the reflection trigger suppressed. On staleness, the LLM re-evaluates the raw environment state without injecting any prior failure context.
\end{enumerate}

\subsubsection{Subtask Progression: Solving the DAG}

\Cref{tab:hrl_breakthrough} presents the exact environment step at which each subtask was first completed across the three configurations.

\begin{table}[H]
\centering
\caption{First breakthrough step (millions of environment steps) for each subtask. \textbf{Bold} indicates the fastest configuration per subtask.}
\label{tab:hrl_breakthrough}
\begin{tabular}{lccc}
\hline
\textbf{Subtask} & \textbf{No-LoRA} & \textbf{QLoRA + Refl.} & \textbf{QLoRA -- No Refl.} \\
\hline
Wood       & \textbf{0.16M} & 0.29M & 0.26M \\
Stone      & \textbf{0.16M} & 0.29M & 0.26M \\
Pickaxe    & 2.10M & \textbf{0.66M} & 0.72M \\
Bridge     & 3.60M & \textbf{1.15M} & 1.18M \\
Iron       & 4.82M & \textbf{2.85M} & 3.01M \\
Sword      & 4.82M & \textbf{2.88M} & 3.15M \\
Armor      & 4.82M & \textbf{2.98M} & 3.15M \\
Enemy      & 4.82M & \textbf{2.88M} & 3.24M \\
Gold       & 11.21M & 22.81M & \textbf{6.29M} \\
\hline
\end{tabular}
\end{table}

Several patterns emerge from \Cref{tab:hrl_breakthrough}. First, all three configurations solve the early-game tasks (Wood, Stone) almost identically, with the No-LoRA base slightly faster due to lower inference overhead. Second, the QLoRA models exhibit a clear advantage in the mid-game: both fine-tuned runs reach Iron approximately 1.8M steps before the base model, demonstrating that fine-tuning improves the LLM's ability to sequence Options coherently through the dependency chain.

The most striking result appears at the Gold subtask. Despite reaching every prior milestone \emph{fastest}, the QLoRA model with reflection enabled takes 22.81M steps to first mine Gold---nearly \textbf{3.6$\times$ slower} than the same model without reflection (6.29M). The No-LoRA base first reaches Gold at 11.21M but ultimately fails to sustain it (final average: 0\%, peak: 8.1\%).

\begin{figure}[H]
    \centering
    \begin{subfigure}{0.32\textwidth}
        \includegraphics[width=\textwidth]{../../plots/hrl_reflection_ablation/hrl_ablation_subtask_Wood.png}
        \caption{Wood}
    \end{subfigure}
    \hfill
    \begin{subfigure}{0.32\textwidth}
        \includegraphics[width=\textwidth]{../../plots/hrl_reflection_ablation/hrl_ablation_subtask_Stone.png}
        \caption{Stone}
    \end{subfigure}
    \hfill
    \begin{subfigure}{0.32\textwidth}
        \includegraphics[width=\textwidth]{../../plots/hrl_reflection_ablation/hrl_ablation_subtask_Pickaxe.png}
        \caption{Pickaxe}
    \end{subfigure}
    \vskip 0.5\baselineskip
    \begin{subfigure}{0.32\textwidth}
        \includegraphics[width=\textwidth]{../../plots/hrl_reflection_ablation/hrl_ablation_subtask_Iron.png}
        \caption{Iron}
    \end{subfigure}
    \hfill
    \begin{subfigure}{0.32\textwidth}
        \includegraphics[width=\textwidth]{../../plots/hrl_reflection_ablation/hrl_ablation_subtask_Sword.png}
        \caption{Sword}
    \end{subfigure}
    \hfill
    \begin{subfigure}{0.32\textwidth}
        \includegraphics[width=\textwidth]{../../plots/hrl_reflection_ablation/hrl_ablation_subtask_Armor.png}
        \caption{Armor}
    \end{subfigure}
    \vskip 0.5\baselineskip
    \begin{subfigure}{0.32\textwidth}
        \includegraphics[width=\textwidth]{../../plots/hrl_reflection_ablation/hrl_ablation_subtask_Bridge.png}
        \caption{Bridge}
    \end{subfigure}
    \hfill
    \begin{subfigure}{0.32\textwidth}
        \includegraphics[width=\textwidth]{../../plots/hrl_reflection_ablation/hrl_ablation_subtask_Enemy.png}
        \caption{Enemy}
    \end{subfigure}
    \hfill
    \begin{subfigure}{0.32\textwidth}
        \includegraphics[width=\textwidth]{../../plots/hrl_reflection_ablation/hrl_ablation_subtask_Gold.png}
        \caption{Gold}
    \end{subfigure}
    \caption{Subtask completion rates across all nine environment milestones for the three HRL configurations. The dependency chain flows left-to-right, top-to-bottom.}
    \label{fig:hrl_all_subtasks}
\end{figure}

\Cref{fig:hrl_all_subtasks} reveals a cascading pattern. In the No-LoRA run, Iron, Sword, Armor, and Enemy all unlock simultaneously at step 4.82M---suggesting the base model issued a single broad Option that happened to cover multiple milestones at once. However, this ``shotgun'' approach fails to sustain Gold extraction: the No-LoRA peak of 8.1\% rapidly collapses to 0\%, indicating the model cannot maintain the precise late-game option sequencing required for the final milestone.

Both QLoRA models achieve 100\% peak Gold completion, but settle at approximately 50\% in the final 100 updates. This ceiling is not a model failure; rather, it reflects the inherent difficulty of mining Gold within the episode's 300-step horizon after completing all prerequisites---approximately half of episodes simply run out of time.

\subsubsection{Reward Trajectories}

\begin{figure}[H]
    \centering
    \includegraphics[width=\textwidth]{../../plots/hrl_reflection_ablation/hrl_ablation_rewards.png}
    \caption{Average Extrinsic Reward across HRL configurations. Both QLoRA runs dramatically outperform the No-LoRA base, which plateaus at 0.030 due to its inability to sustain Gold mining.}
    \label{fig:hrl_rewards}
\end{figure}

The reward trajectories (\Cref{fig:hrl_rewards}) quantify the performance gap. While the Dynamic Shaping architectures (Part~1) plateaued at an average reward of 0.076, the QLoRA HRL runs achieve peak rewards of 0.335 (No Reflection) and 0.330 (With Reflection)---a \textbf{4.4$\times$ improvement}. The No-LoRA base plateaus at just 0.030, confirming that while the base model can generate syntactically valid Options, it lacks the fine-tuned coherence to sustain late-game performance.

Crucially, the No-Reflection run reaches its peak reward at step 28.4M, while the With-Reflection run does not peak until step 59.9M---more than twice as late. This 2.1$\times$ sample efficiency gap is the first quantitative signal that the reflection mechanism is actively harmful.

\subsubsection{Stabilizing Non-Stationary Objectives via KL-Divergence Guarding}
\label{sec:kl_guard}

In the HRL architecture, the transition between discrete macro-actions (Options) induces abrupt shifts in the intrinsic reward landscape. From the perspective of the low-level PPO workers, the objective function is highly non-stationary. To prevent catastrophic forgetting of previously learned kinetic policies, we bound the surrogate objective update using a strict Kullback-Leibler (KL) divergence guard.

During the optimization epochs, we continuously monitor the approximate KL divergence between the behavioral policy $\pi_{\theta_{\text{old}}}$ and the updating policy $\pi_{\theta}$:
$$D_{\text{KL}}^{\text{approx}} = \mathbb{E}\left[ -\log \left( \frac{\pi_{\theta}(a|s, o)}{\pi_{\theta_{\text{old}}}(a|s, o)} \right) \right]$$

If $D_{\text{KL}}^{\text{approx}}$ exceeds the predefined trust region threshold ($\delta = 0.015$), the gradient updates for the current batch are immediately aborted.

\begin{figure}[H]
    \centering
    \includegraphics[width=\textwidth]{../../plots/hrl_reflection_ablation/hrl_ablation_kl_spikes.png}
    \caption{Maximum approximate KL divergence per update across the three HRL configurations. The dashed red line marks the early-stopping threshold $\delta = 0.015$.}
    \label{fig:kl_spikes}
\end{figure}

\begin{table}[H]
\centering
\caption{KL divergence statistics and policy stability metrics across HRL configurations.}
\label{tab:kl_stats}
\begin{tabular}{lccc}
\hline
\textbf{Metric} & \textbf{No-LoRA} & \textbf{QLoRA + Refl.} & \textbf{QLoRA -- No Refl.} \\
\hline
Mean $\max D_{\text{KL}}$       & 0.0063 & 0.0063 & 0.0073 \\
Peak $\max D_{\text{KL}}$       & 0.019  & 0.021  & \textbf{0.107} \\
KL early-stop triggers           & 11 / 2{,}000 & 21 / 2{,}000 & \textbf{152 / 2{,}000} \\
Staleness resets (total)         & 21{,}492 & \textbf{54{,}534} & 39{,}137 \\
LLM queries dispatched           & 6{,}710 & 9{,}160 & \textbf{13{,}710} \\
Final TD error (last 100 updates)& 0.047  & 0.039  & \textbf{0.023} \\
\hline
\end{tabular}
\end{table}

\Cref{fig:kl_spikes} and \Cref{tab:kl_stats} reveal a counter-intuitive relationship between KL instability and final performance. The No-Reflection run triggers the KL early-stopping guardrail \textbf{7.2$\times$ more often} than the With-Reflection run (152 vs.\ 21 triggers), yet achieves the lowest final TD error (0.023 vs.\ 0.039) and the highest reward. Conversely, the With-Reflection run appears superficially ``stable'' (low KL, few early stops) but accumulates \textbf{2.5$\times$ more staleness resets} (54{,}534 vs.\ 39{,}137).

This paradox is explained by a fundamental difference in how each configuration handles option transitions. Without reflection, the LLM commits to long-horizon Options and rarely switches. When it does switch---for example, from \texttt{CRAFT\_SWORD} to \texttt{MINE\_GOLD}---the policy shift is large, producing the KL spikes visible in \Cref{fig:kl_spikes} (peak: 0.107, or $7\times$ the threshold). But the guardrail immediately catches and bounds the update. The policy absorbs the new objective gradually over subsequent updates. This pattern of \emph{infrequent, large, guarded} transitions produces stable long-term learning.

With reflection enabled, the mechanism creates a destructive feedback loop: staleness triggers a reflection prompt $\rightarrow$ the LLM switches Options $\rightarrow$ agents receive conflicting intrinsic rewards $\rightarrow$ the new Option goes stale $\rightarrow$ another reflection fires. Each individual switch is small enough (peak KL = 0.021) to slip beneath the guardrail threshold, but the rapid succession of switches prevents the policy from ever converging on a stable objective. This \emph{death by a thousand cuts} is invisible to the KL guard yet devastates sample efficiency.

\subsubsection{The Reflection Contradiction}
\label{sec:reflection_ablation}

The staleness data (\Cref{tab:kl_stats}) provides the mechanistic explanation for the performance gap. The With-Reflection run dispatches fewer LLM queries (9{,}160 vs.\ 13{,}710) yet accumulates 39\% more staleness resets. This means each reflection-triggered query produces an Option that fails faster---the LLM is actively generating worse plans when told it previously failed.

The root cause lies in how the Oracle training dataset was constructed. During QLoRA fine-tuning, the dataset synthetically injected counterfactual failures: the training prompt stated \emph{``You previously chose [Incorrect Option]. It failed.''} The model then learned to output the correct Option. Consequently, the LLM associated the reflection trigger with a strong prior: \textbf{my previous choice was logically wrong, and I must switch to a different option.}

During inference, however, option staleness is frequently triggered not by a logical error, but by kinetic delays---agents navigating around the river terrain, queuing at resource nodes, or simply needing more steps to execute a fundamentally correct plan. When the orchestrator fires the reflection prompt (\emph{``You previously chose \texttt{MINE\_GOLD}. After 150 steps, it failed.''}), the LLM, conditioned by its training data, assumes it made a reasoning error. It abandons the correct long-horizon option and switches to an alternative, derailing the agents' execution.

We term this failure mode the \emph{Reflection Contradiction}: the architecture inadvertently trained the LLM to doubt its own correct logic. Disabling the reflection mechanism allows the LLM to maintain stable plans---on re-query, it re-confirms its current Option rather than compulsively switching.

\begin{figure}[H]
    \centering
    \includegraphics[width=\textwidth]{../../plots/hrl_reflection_ablation/hrl_ablation_td_error.png}
    \caption{Temporal Difference Error Variance. The With-Reflection configuration exhibits persistent high variance from contradictory option switches, while the No-Reflection run converges to a stable 0.023.}
    \label{fig:hrl_td_error}
\end{figure}

The downstream effect is quantified in the TD Error Variance (\Cref{fig:hrl_td_error}). Each contradictory option switch reshapes the intrinsic reward landscape, forcing the Critic network to chase a non-stationary target. The With-Reflection configuration maintains a mean TD error of 0.064 throughout training---34\% higher than the No-Reflection run (0.048). By the final 100 updates, this gap widens further: the No-Reflection Critic converges to 0.023, while the With-Reflection Critic remains at 0.039, indicating persistent instability even at the end of training.

Ultimately, the HRL architecture with QLoRA fine-tuning but \emph{without} counterfactual reflection represents the strongest configuration: it solves the complete dependency chain 2.1$\times$ faster, achieves the lowest policy variance, and reaches peak reward with half the training budget. The KL-divergence guardrail proves essential for absorbing the inherent non-stationarity of option transitions, while the reflection mechanism---despite its theoretical appeal---actively undermines long-horizon planning by injecting ungrounded self-doubt into the Meta-Controller's reasoning.

% ==========================================================================
\subsection{Part 3: Robustness and Semantic Adaptation}
\label{sec:robustness_and_adaptation}

To validate the structural components of the Mixing Board architecture, we isolate two critical variables: the numerical stability provided by temporal smoothing, and the zero-shot semantic adaptation capabilities of the LLM Meta-Controller.

\subsubsection{Ablation: The Necessity of EMA Smoothing}
\label{sec:ema_ablation}

In Chapter 4, we posited that feeding raw, discrete priority logits directly into the potential-based shaping function would induce severe non-stationarity. To empirically verify this, an ablation run was conducted wherein the Exponential Moving Average (EMA) layer was disabled.

\begin{figure}[H]
    \centering
    \includegraphics[width=0.48\textwidth]{../../plots/llm_dynamic_ablations_and_shifts/ablation_rewards.png}
    \hfill
    \includegraphics[width=0.48\textwidth]{../../plots/llm_dynamic_ablations_and_shifts/ablation_td_error.png}
    \caption{Comparison of training stability with and without EMA smoothing ($\alpha=0.2$). Disabling smoothing induces a catastrophic spike in Temporal Difference (TD) error variance (right).}
    \label{fig:ema_ablation_td}
\end{figure}

\begin{figure}[H]
    \centering
    \includegraphics[width=0.48\textwidth]{../../plots/llm_dynamic_ablations_and_shifts/ablation_subtask_Iron.png}
    \hfill
    \includegraphics[width=0.48\textwidth]{../../plots/llm_dynamic_ablations_and_shifts/ablation_subtask_Sword.png}
    \caption{Subtask completion rates for Iron (left) and Sword (right). The variance spike directly correlates with catastrophic forgetting in advanced subtasks.}
    \label{fig:ema_ablation_subtasks}
\end{figure}

As shown in Figure~\ref{fig:ema_ablation_td}, removing the EMA filter results in a catastrophic spike in TD-error variance. The ablation itself was executed by passing a \texttt{skip\_ema=True} flag to the \texttt{LLMOrchestrator}, which bypassed the exponential smoothing step entirely:

\begin{lstlisting}[style=code, language=Python, caption={EMA Ablation implementation in \texttt{llm/orchestrator.py}.}, label={lst:ema_ablation}]
# EMA Smoothing (alpha = 0.2) -- skipped in ablation mode
if self.skip_ema:
    self.active_weights = parsed_weights
else:
    alpha = 0.2
    for k in self.active_weights:
        self.active_weights[k] = alpha * parsed_weights[k] + (1 - alpha) * self.active_weights[k]
\end{lstlisting}

The Critic network struggles to track a reward objective that shifts instantaneously. Crucially, Figure~\ref{fig:ema_ablation_subtasks} illustrates the behavioral consequence: the non-stationarity causes the agents to temporarily forget complex subtask sequences, dropping their `Iron` and `Sword` completion rates significantly in the mid-stages of training. However, it is worth noting that the No-EMA agents do eventually recover and climb back to a 100\% success rate on these tasks. Thus, while the agents can mathematically overcome the noise, this proves that the EMA smoothing equation is highly helpful for bridging discrete LLM outputs with continuous RL gradients, enabling faster convergence and preventing catastrophic forgetting.

\begin{figure}[H]
    \centering
    \includegraphics[width=0.8\textwidth]{../../plots/llm_dynamic_ablations_and_shifts/ablation_all_subtasks_bar.png}
    \caption{Maximum subtask completion rates achieved across all runs. Despite the early variance spikes, No-EMA eventually masters all tasks identical to the baseline.}
    \label{fig:ema_ablation_all_subtasks}
\end{figure}

\subsubsection{Environment Shift: The Paradox of Semantic Over-Confidence}
\label{sec:semantic_adaptation}

A fundamental critique of using LLMs in MARL is whether they rely on genuine semantic world-models or merely overfit to node names. To test this, the environment was silently modified at Update 500: the `Workbench` zone was renamed to `Furnace` in the LLM's prompt space. A secondary control run replaced `Workbench` with a nonsense string, `Zylorg`.

This shift was executed by dynamically injecting an alias mapping into the \texttt{LLMOrchestrator} mid-training, forcing the LLM to output weights for the alias rather than the canonical node name, which the orchestrator then reversed and fed to the PPO agents:

\begin{lstlisting}[style=code, language=Python, caption={Dynamic semantic shift implementation in \texttt{llm/orchestrator.py}.}, label={lst:env_shift_code}]
# Build reverse alias map: e.g. "w_furnace" -> "w_workbench"
alias_reverse = {}
for canonical, alias in self.zone_aliases.items():
    alias_reverse[f"w_{alias.lower()}"] = f"w_{canonical.lower()}"

# Extract weights: try canonical key first, then fallback to aliased key
alias_k = next((ak for ak, ck in alias_reverse.items() if ck == k), None)
val = raw.get(k, raw.get(alias_k))
\end{lstlisting}

The initial hypothesis assumed the LLM would cleanly map "Furnace" to crafting success, while failing on "Zylorg". However, empirical results revealed a counter-intuitive phenomenon: \textbf{semantic over-confidence induced severe reward hacking, while semantic ignorance acted as a natural regularizer.}

\begin{figure}[H]
    \centering
    \includegraphics[width=0.48\textwidth]{../../plots/llm_dynamic_ablations_and_shifts/shift_subtask_Iron.png}
    \hfill
    \includegraphics[width=0.48\textwidth]{../../plots/llm_dynamic_ablations_and_shifts/shift_subtask_Sword.png}
    \caption{Subtask completion rates during semantic shifts. The `Furnace` shift traps agents in a local optimum, halting `Iron` and `Sword` progression, while the `Zylorg` shift organically recovers.}
    \label{fig:semantic_shift_subtasks}
\end{figure}

\begin{figure}[H]
    \centering
    \includegraphics[width=0.8\textwidth]{../../plots/llm_dynamic_ablations_and_shifts/shift_all_subtasks_bar.png}
    \caption{Maximum subtask completion rates for the Shift group. Furnace masters initial tasks (Wood, Stone, Pickaxe) but completely fails advanced tasks requiring exploration (Iron, Sword).}
    \label{fig:shift_all_subtasks}
\end{figure}

As seen in Figures~\ref{fig:semantic_shift_subtasks} and~\ref{fig:shift_all_subtasks}, the `Zylorg` control unexpectedly succeeded across almost all subtasks, achieving near-parity with the baseline's extrinsic reward. When the LLM encountered "Zylorg", it lacked a semantic prior. Consequently, it assigned it a balanced priority weight alongside other nodes, focusing its reasoning on recognizable end-goals like enemies. 

\begin{lstlisting}[style=llmprompt, caption={Actual JSONL log excerpt during the 'Zylorg' environment shift.}, label={lst:zylorg_log}]
"reasoning": "Enemy defeat is the current bottleneck with no progress.",
"w_wood": 0.8,
"w_stone": 0.8,
"w_zylorg": 0.9,
"w_iron": 0.9,
"w_enemy": 1.2
\end{lstlisting}

Because \texttt{w\_zylorg} was not artificially inflated, the MAPPO agents, driven by the baseline extrinsic rewards for crafting a Sword, naturally explored the environment, discovered the `Zylorg` station was the required prerequisite, and solved the bottleneck via standard PPO exploration. The LLM's ignorance prevented it from interfering with the RL algorithm's foundational mechanics.

Conversely, the `Furnace` shift failed catastrophically on advanced tasks (`Iron`, `Sword`) due to reward hacking. The LLM correctly deduced zero-shot that a Furnace serves an identical semantic purpose to a crafting bench. It demonstrated this reasoning in the raw intervention logs:

\begin{lstlisting}[style=llmprompt, caption={Actual JSONL log excerpt during the 'Furnace' environment shift.}, label={lst:furnace_log}]
"reasoning": "Agents are not crafting Pickaxe or Sword, indicating a bottleneck at the Workbench.",
"w_furnace": 0.4,
"w_iron": 0.2
\end{lstlisting}

However, by mapping the Furnace to the crafting bottleneck, the LLM heavily skewed the Potential-Based Reward Shaping (PBRS) landscape. The PPO agents successfully crafted the Pickaxe (Figure~\ref{fig:shift_all_subtasks}), but subsequently discovered they could accumulate massive intrinsic returns simply by navigating back to the Furnace and idling there. They optimized for the dense shaping signal at the expense of the sparse extrinsic reward, halting all exploration for `Iron` and `Sword`. 

This finding highlights a critical vulnerability in neuro-symbolic MARL architectures: while LLMs possess exceptional zero-shot semantic mapping, coupling high-confidence oracle weights with dense gradient shaping can inadvertently trap low-level agents in local optima.


